\documentclass[11pt]{article}
\usepackage{amsmath}
\usepackage{amssymb}
\usepackage{mathtools}
\usepackage{geometry}
\usepackage{booktabs}
\geometry{margin=1in}

\title{Temporal LDA: Density-Weighted Topic Model}
\author{}
\date{}

\begin{document}

\maketitle

\begin{abstract}
Standard topic models treat all documents equally, causing topics from sparse temporal periods to be overwhelmed by dense periods.
We present a density-weighted variant of Latent Dirichlet Allocation (\textbf{Temporal LDA}) that uses Gaussian kernel density estimation on word mass to compute an inverse-density weight for each document.
The key design is asymmetric: document-topic counts remain unweighted (preserving realistic document composition), while topic-word counts are weighted (equalizing the contribution of sparse and dense periods to learned vocabularies).
Using collapsed Gibbs sampling on synthetic corpora with 1:20 document imbalance, we show that weighting recovers all six ground-truth topics where standard LDA collapses the sparse era into a single blob.
With fewer topics than ground truth ($K=4$), the weighted model splits its budget symmetrically (2~sparse + 2~dense) while standard LDA devotes 3 of 4~topics to the dense era.
\end{abstract}

%% ===================================================================
\section{Problem: Temporal Imbalance in Topic Models}
%% ===================================================================

In historical document collections the number of available sources varies dramatically across time.
A corpus might contain 50 documents from a sparsely documented era and 500 from a densely documented one---a 1:10 imbalance in both document count and total word mass.

Standard LDA maximises the joint likelihood over all tokens equally:
\[
\log P(\mathbf w) = \sum_{d=1}^{D} \sum_{n=1}^{N_d} \log P(w_{d,n} \mid \theta_d, \boldsymbol\beta)
\]
Because the 800-document period contributes ten times as many tokens, the model learns topic-word distributions that explain the dense period well while collapsing or ignoring themes unique to the sparse period.

%% ===================================================================
\section{Density Weighting}
%% ===================================================================

We weight each document inversely by the local temporal density of word mass around its year.

\paragraph{Gaussian kernel.}
Define a centred Gaussian kernel with bandwidth~$\sigma$ (in years):
\[
K(t) = \exp\!\left( -\frac{t^2}{2\sigma^2} \right)
\]

\paragraph{Word-mass density.}
For each document~$d$ at year~$t_d$ with word count~$\ell_d$:
\[
\rho(t_d)
  = \frac{\displaystyle\sum_{d'=1}^{D} \ell_{d'}\, K(t_{d'} - t_d)}
         {\displaystyle\left(\sum_{d'=1}^{D} \ell_{d'}\right)
          \cdot \sum_{t=y_{\min}}^{y_{\max}} K(t - t_d)}
\]

The \textbf{density weight} is the inverse of~$\rho$, normalised so the mean weight across documents equals~1:
\[
\widetilde w_d = \frac{1}{\rho(t_d)},
\qquad
w_d = \frac{\widetilde w_d}{\frac{1}{D}\sum_{d'}\widetilde w_{d'}}
\]

On the 50/500 test data with $\sigma=10$, sparse-era documents receive $w_d \approx 7.3$ and dense-era documents receive $w_d \approx 0.37$.

%% ===================================================================
\section{Generative Model}
%% ===================================================================

The generative model is LDA with a symmetric topic-word prior and an asymmetric, time-dependent document-topic prior:
\[
\begin{aligned}
\boldsymbol\beta_k &\sim \operatorname{Dir}(\beta) &&\text{for } k = 1,\ldots,K \\
\boldsymbol\theta_d &\sim \operatorname{Dir}\!\bigl(\alpha_1(y_d),\ldots,\alpha_K(y_d)\bigr) &&\text{for } d = 1,\ldots,D \\
z_{d,n} &\sim \operatorname{Cat}(\boldsymbol\theta_d) \\
w_{d,n} &\sim \operatorname{Cat}(\boldsymbol\beta_{z_{d,n}})
\end{aligned}
\]
where $\boldsymbol\alpha(y) = (\alpha_1(y),\ldots,\alpha_K(y))$ is a per-topic document-topic prior that varies with year~$y$, and $\beta$ is a scalar topic-word prior.
The per-topic structure lets the model learn which topics are active in each era (inactive topics receive $\alpha_k \approx 0$), while the temporal dependence lets the concentration vary across time.
The density weight enters only during inference, not in the generative process.

%% ===================================================================
\section{Collapsed Gibbs Sampling with Asymmetric Weighting}
%% ===================================================================

We use collapsed Gibbs sampling, integrating out~$\boldsymbol\theta$ and~$\boldsymbol\beta$ to work directly with sufficient-statistic counts.

\subsection{Asymmetric Count Bookkeeping}

Three count arrays are maintained:
\begin{itemize}
\item $n_{d,k}$\;---\;tokens in document~$d$ assigned to topic~$k$ \textbf{(unweighted, $\pm\,1$)}
\item $n_{k,v}$\;---\;occurrences of word~$v$ in topic~$k$ \textbf{(weighted, $\pm\,w_d$)}
\item $n_{k}$\;---\;total tokens in topic~$k$ \textbf{(weighted, $\pm\,w_d$)}
\end{itemize}
When a token is removed or added during sampling, $n_{d,k}$ changes by~$1$ while $n_{k,v}$ and~$n_k$ change by~$w_d$.

\subsection{Sampling Equation}

For each token~$(d,n)$ with current assignment~$z_{d,n}$:

\begin{enumerate}
\item \textbf{Remove} the token from all three count arrays.
\item \textbf{Sample} a new topic from
\[
P(z_{d,n} = k \mid \mathbf z^{-(d,n)}, \mathbf w)
  \;\propto\;
  \bigl(n_{d,k}^{-(d,n)} + \alpha_k(y_d)\bigr)
  \;\cdot\;
  \frac{n_{k,v}^{-(d,n)} + \beta}
       {n_{k}^{-(d,n)} + \beta\,V}
\]
where $V$ is the vocabulary size and $\alpha_k(y_d)$ is the per-topic prior for the document's year.
\item \textbf{Add} the token back under the newly sampled topic.
\end{enumerate}

\subsection{Why Asymmetric Weighting}

\begin{enumerate}
\item \textbf{Preserves document composition.}
Unweighted~$n_{d,k}$ means each document's topic mixture~$\theta_d$ reflects its actual word content, not an inflated count.

\item \textbf{Balances vocabulary learning.}
Weighted~$n_{k,v}$ gives sparse-period documents a larger per-token influence on learned topic-word distributions, counteracting the raw imbalance in token counts.
\end{enumerate}

%% ===================================================================
\section{Bayesian Interpretation}
%% ===================================================================

The weighting scheme has a precise interpretation in terms of the posterior over topic-word distributions.
This section makes the connection explicit.

\subsection{Standard Collapsed Gibbs and Sufficient Statistics}

In standard collapsed LDA the topic-word distributions~$\boldsymbol\beta_k$ are integrated out analytically.
Conditioned on topic assignments~$\mathbf z$, the posterior over~$\boldsymbol\beta_k$ is Dirichlet:
\[
\boldsymbol\beta_k \mid \mathbf z, \mathbf w
  \;\sim\;
  \operatorname{Dir}\!\bigl(\beta + n_{k,1},\;\ldots,\;\beta + n_{k,V}\bigr)
\]
where $n_{k,v} = \#\{(d,n) : z_{d,n}=k,\; w_{d,n}=v\}$ is the unweighted count of word~$v$ assigned to topic~$k$.
Every token contributes equally: one unit to the count of its word in its topic.

\subsection{What Weighting Changes}

In our scheme a token from document~$d$ contributes $w_d$ instead of~$1$ to the topic-word counts:
\[
\tilde n_{k,v}
  = \sum_{\substack{(d,n):\\z_{d,n}=k,\;w_{d,n}=v}} w_d
\]

If we substitute~$\tilde n_{k,v}$ into the Dirichlet posterior we get
\[
\boldsymbol\beta_k \mid \mathbf z, \mathbf w
  \;\sim\;
  \operatorname{Dir}\!\bigl(\beta + \tilde n_{k,1},\;\ldots,\;\beta + \tilde n_{k,V}\bigr)
\]
Dirichlet parameters need not be integers, so this remains a well-defined distribution.
The point estimate used in the code (the posterior mean) is
\[
\hat\beta_{k,v}
  = \frac{\tilde n_{k,v} + \beta}{\tilde n_k + \beta\,V}
\]
which is exactly the formula in the output routine.

\subsection{Interpretation: Fractional Observations}

The Dirichlet-Multinomial conjugate update treats each count unit as one observation of a word.
By adding~$w_d$ instead of~$1$, we are telling the model that a token from a sparse period is worth~$w_d$ observations (e.g.\ $7.3$ for sparse-era documents) while a token from a dense period is worth only~$w_d \approx 0.37$ observations.

This is equivalent to \textbf{importance weighting} on the sufficient statistics: each real token is reweighted so that the effective corpus seen by the topic-word posterior is temporally balanced.

\subsection{What Is \emph{Not} Reweighted}

Document-topic counts~$n_{d,k}$ remain unweighted.
The posterior over each document's topic mixture is therefore
\[
\boldsymbol\theta_d \mid \mathbf z
  \;\sim\;
  \operatorname{Dir}\!\bigl(\alpha_1(y_d) + n_{d,1},\;\ldots,\;\alpha_K(y_d) + n_{d,K}\bigr)
\]
with the true integer counts and the per-topic prior for the document's year.
This means $\theta_d$ reflects the document's actual topic proportions, not an inflated version.
The asymmetry is deliberate: we want to correct the \emph{global} vocabulary posteriors for temporal imbalance while leaving \emph{per-document} composition estimates truthful.

\subsection{Effect on the Gibbs Conditional}

The Gibbs sampling equation
\[
P(z_{d,n} = k) \;\propto\;
  \underbrace{\bigl(n_{d,k}^{-(d,n)} + \alpha_k(y_d)\bigr)}_{\text{unweighted}}
  \;\cdot\;
  \underbrace{\frac{\tilde n_{k,v}^{-(d,n)} + \beta}{\tilde n_k^{-(d,n)} + \beta\,V}}_{\text{weighted}}
\]
involves the product of an unweighted document-topic term and a weighted topic-word term.
The document-topic term pulls the token toward topics already present in the document (a local signal).
The topic-word term evaluates how well each topic explains this particular word given the rebalanced global vocabulary---a global signal in which sparse-period words carry amplified influence.

During burn-in this drives the sampler toward configurations where topic-word distributions are shaped equally by both periods, while each document's internal mixture remains faithful to its own content.

%% ===================================================================
\section{What the Weighting Equalises (and What It Does Not)}
%% ===================================================================

\subsection{Claim: Uniform Effective Word Mass}

The weighting has a single provable consequence:
\emph{the weighted word mass per unit time is approximately constant across the timeline.}
That is, every year exerts equal influence on the topic-word sufficient statistics, regardless of how many documents survive from that year.

\subsection{Sketch of Proof}

Define the \textbf{unnormalised word-mass density} at continuous time~$y$:
\[
\hat\rho(y) = \sum_{d=1}^{D} \ell_d \, K_\sigma(y - y_d)
\]
This is the kernel-smoothed total word mass at~$y$.
The weight of each document is $w_d \propto 1/\hat\rho(y_d)$ (with a subsequent normalisation so that $\bar w = 1$).

Now consider the \textbf{weighted word-mass density}---the effective word mass at~$y$ after reweighting:
\[
\hat\rho_w(y)
  = \sum_{d=1}^{D} \ell_d \, w_d \, K_\sigma(y - y_d)
  \;\propto\; \sum_{d=1}^{D} \frac{\ell_d}{\hat\rho(y_d)}\, K_\sigma(y - y_d)
\]

Because $\hat\rho$ is itself a kernel smooth with bandwidth~$\sigma$, it varies slowly on the scale of~$\sigma$.
For the documents that contribute non-negligibly to the sum at~$y$ (those with $|y_d - y| \lesssim \sigma$), we have $\hat\rho(y_d) \approx \hat\rho(y)$, so
\[
\hat\rho_w(y)
  \;\approx\; \frac{1}{\hat\rho(y)} \sum_{d=1}^{D} \ell_d \, K_\sigma(y - y_d)
  \;=\; \frac{\hat\rho(y)}{\hat\rho(y)}
  \;=\; 1
\]

The approximation is exact in the limit $\sigma \to \infty$ (where $\hat\rho$ becomes constant) and remains tight whenever the density varies smoothly relative to~$\sigma$.
In practice the Gaussian kernel's rapid decay makes the approximation excellent for any corpus where the year distribution does not have sharp discontinuities narrower than~$\sigma$.

\subsection{Consequence for Topic Estimation}

After the sampler converges, the weighted topic-word counts are
\[
\tilde n_{k,v}
  = \sum_{\substack{(d,n):\\z_{d,n}=k,\;w_{d,n}=v}} w_d
\]
The total weighted count that topic~$k$ accumulates from documents near year~$y$ is proportional to the weighted word mass from that year multiplied by the fraction of that mass assigned to~$k$.
Because $\hat\rho_w(y) \approx 1$, \textbf{each year contributes approximately equal effective word mass to every topic-word posterior},
with the cross-year variation determined only by how much each year actually uses topic~$k$, not by how many documents happen to survive from that year.

\subsection{What This Does \emph{Not} Guarantee}

The weighting equalises \emph{influence}, not \emph{fit}.
It does not follow that the model explains every word equally well across time.
If the true vocabulary genuinely differs between periods---which it does; that is the entire reason distinct topics exist---then the per-token log-likelihood will still vary by era.
Two kinds of variation remain:

\begin{enumerate}
\item \textbf{Topical variation.}
  A word that belongs to a well-separated era-specific topic will be explained well in its own period and poorly outside it.
  This is correct behaviour, not a defect.

\item \textbf{Model capacity.}
  If $K$ is too small to represent all true themes, the model must merge topics.
  Which topics get merged depends on the data, and weighting ensures the merge is not biased toward the dense period---but some loss of fit is inevitable wherever merging occurs.
\end{enumerate}

The precise guarantee is: \emph{the topics are estimated as if the corpus had uniform temporal coverage.}
Dense periods cannot dominate the topic-word distributions; sparse periods cannot be drowned out.
The compromise the model makes when fitting a single set of~$K$ topics is balanced across time, rather than skewed toward whichever era contributes more tokens.

%% ===================================================================
\section{Implementation}
%% ===================================================================

\subsection{Input Files}

\textbf{documents.txt}\;---\;one line per document:
\begin{verbatim}
year  length  word_id_1  word_id_2  ...
\end{verbatim}

\textbf{vocab.txt}\;---\;one word per line, ordered by word ID.

\textbf{metadata.txt}\;---\;corpus facts (read by the program):
\begin{verbatim}
num_documents=550
vocab_size=18
year_min=1850
year_max=1950
\end{verbatim}

\subsection{Compilation and CLI}

\begin{verbatim}
make

./temporal_lda --docs documents.txt \
               --vocab vocab.txt    \
               --metadata metadata.txt \
               --output results/    \
               --K 4 --iterations 1000 --seed 99 \
               --alpha 1.0 --beta 1.0 --sigma 50.0 \
               --optimize-interval 50 --converge 1e-6
\end{verbatim}

Pass \texttt{--no-weighting} to run standard (unweighted) LDA.
Pass \texttt{--converge $\epsilon$} to stop early once the log-likelihood stabilises (\texttt{--iterations} becomes the upper bound).
Pass \texttt{--optimize-interval $N$} to re-estimate $\alpha$ and $\beta$ every $N$~Gibbs iterations using Minka's fixed-point update (see below).
Pass \texttt{--local-alpha FILE} to enable time-dependent $\boldsymbol\alpha(y)$ and write the per-year vectors to \texttt{FILE}; without this flag $\boldsymbol\alpha$ is a single global vector (the same asymmetric prior at every year).
All model parameters ($K$, $\alpha$, $\beta$, $\sigma$, seed, iterations) are set via CLI flags; the metadata file supplies only corpus dimensions.

\subsection{Hyperparameter Optimisation}

When \texttt{--optimize-interval $N$} is set, the sampler periodically re-estimates $\boldsymbol\alpha$ and $\beta$ using Minka's fixed-point iteration for asymmetric Dirichlet priors.

\paragraph{Asymmetric Minka update.}
The sampler maintains a $K$-vector $\boldsymbol\alpha(y) = (\alpha_1(y),\ldots,\alpha_K(y))$ for each year~$y$.
The update uses a shared denominator across all topics and a per-topic numerator:
\[
\alpha_k(y)^{\text{new}}
  = \alpha_k(y)\;
    \frac{\displaystyle\sum_{d=1}^{D} K_\sigma(t_d - y)\,
          \bigl[\psi(n_{d,k}+\alpha_k(y))-\psi(\alpha_k(y))\bigr]}
         {\displaystyle\sum_{d=1}^{D} K_\sigma(t_d - y)\,
          \bigl[\psi(N_d+\alpha_{\Sigma}(y))-\psi(\alpha_{\Sigma}(y))\bigr]}
\]
where $\alpha_{\Sigma}(y) = \sum_{k=1}^K \alpha_k(y)$ is the concentration,
$N_d = \sum_k n_{d,k}$ is the document length,
and $K_\sigma(t) = \exp(-t^2/2\sigma^2)$ is the same Gaussian kernel used for density weighting.

The per-topic structure lets the estimator push inactive topics toward $\alpha_k \approx 0$ while active topics converge to their true Dirichlet concentration.
The shared denominator couples the topics through the total document length, enforcing the constraint that the concentration sum~$\alpha_\Sigma$ must explain the overall sparsity of document-topic mixtures.

\paragraph{Local vs.\ global mode.}
When \texttt{--local-alpha} is passed, each year~$y$ maintains its own $K$-vector~$\boldsymbol\alpha(y)$, kernel-weighted by temporal proximity as above.
Without \texttt{--local-alpha}, a single global $K$-vector is estimated (using all documents with flat weighting) and broadcast to all years.

The kernel bandwidth~$\sigma$ controls the locality of the local estimate:
\begin{itemize}
\item $\sigma \to \infty$: the kernel is flat, recovering the global mode.
\item $\sigma \to 0$: each year's $\boldsymbol\alpha$ depends only on documents from that exact year.
\end{itemize}

\paragraph{Global $\beta$.}
The topic-word prior $\beta$ is updated with the standard (global) Minka iteration:
\[
\beta^{\text{new}}
  = \beta \;\frac{\displaystyle\sum_{k=1}^{K}\sum_{v=1}^{V}
      \bigl[\psi(\tilde n_{k,v}+\beta)-\psi(\beta)\bigr]}
     {\displaystyle V\sum_{k=1}^{K}
      \bigl[\psi(\tilde n_{k}+V\beta)-\psi(V\beta)\bigr]}
\]
using the weighted topic-word counts $\tilde n_{k,v}$ and totals $\tilde n_k$.

Five inner fixed-point steps are applied per optimisation round; a 50-iteration burn-in precedes the first round.
In practice $\alpha(y)$ and $\beta$ converge from their initial values (typically~1.0) to stable values within ${\sim}200$ Gibbs iterations.

\subsection{Convergence Detection}

The \texttt{--converge $\epsilon$} flag enables window-mean convergence detection.
The sampler maintains two consecutive windows of 10~log-likelihood samples, each spanning 100~Gibbs iterations.
After both windows are full, convergence is declared when
\[
\frac{\lvert\bar L_{\text{curr}} - \bar L_{\text{prev}}\rvert}{\lvert\bar L_{\text{curr}}\rvert}
  < \epsilon
\]
where $\bar L$ denotes the window mean.
This smooths out MCMC noise while remaining sensitive to genuine drift in the likelihood surface.

\subsection{Output Files}

\begin{itemize}
\item \texttt{beta.txt}\;---\;topic-word distributions: \texttt{topic word prob} per line.
\item \texttt{theta.txt}\;---\;document-topic distributions: \texttt{doc topic prob} per line.
\item \texttt{weights.txt}\;---\;per-document weights: \texttt{doc year length weight} per line.
\item \texttt{alpha.txt}\;---\;per-year document-topic prior: \texttt{year $\alpha_1$ $\alpha_2$ \ldots\ $\alpha_K$} per line
  ($K$~columns; only with \texttt{--local-alpha}; filename is user-specified).
\end{itemize}

%% ===================================================================
\appendix
\section{Validation Experiment}
%% ===================================================================

We validate on a synthetic corpus generated by \texttt{gen\_test.rb}.
The corpus has two time points---a sparse era and a dense era---with a 1:20 imbalance in total word mass, heterogeneous document-topic sparsity, and six ground-truth topics whose vocabularies are completely disjoint.

\subsection{Data}

Six true topics, each defined by three exclusive words from a vocabulary of 18~symbols:

\begin{center}
\begin{tabular}{@{}l l l@{}}
\toprule
\textbf{Topic} & \textbf{Words} & \textbf{Era} \\
\midrule
T1 & \texttt{a b c} & 1850 (sparse) \\
T2 & \texttt{d e f} & 1850 (sparse) \\
T3 & \texttt{g h i} & 1850 (sparse) \\
T4 & \texttt{A B C} & 1950 (dense)  \\
T5 & \texttt{D E F} & 1950 (dense)  \\
T6 & \texttt{G H I} & 1950 (dense)  \\
\bottomrule
\end{tabular}
\end{center}

\begin{itemize}
\item \textbf{1850 (sparse):} 50~documents of length 100.  Each document draws a topic mixture from $\operatorname{Dir}(0.1, 0.1, 0.1)$ over T1, T2, T3.
\item \textbf{1950 (dense):} 500~documents of length 200.  Each document draws a topic mixture from $\operatorname{Dir}(0.3, 0.3, 0.3)$ over T4, T5, T6.
\end{itemize}

Total: 550~documents, 105{,}000~tokens.
The dense era contributes 20$\times$ as many tokens as the sparse era (100{,}000 vs.\ 5{,}000).
The two eras also have genuinely different document-topic structure: the sparse era's $\operatorname{Dir}(0.1)$ concentrates each document on one or two topics, while the dense era's $\operatorname{Dir}(0.3)$ produces more diffuse mixtures.

\subsection{Configuration}

Both conditions use $K=6$, initial $\alpha=1$, initial $\beta=1$, up to 2{,}000~iterations, \texttt{--optimize-interval\;50} (asymmetric Minka updates every 50~iterations), and \texttt{--converge\;1e-6}.
The weighted run additionally passes \texttt{--local-alpha} to enable per-year asymmetric $\boldsymbol\alpha(y)$ estimation, with $\sigma=10$ (years), giving sparse-era documents a weight of ${\approx}7.3$ and dense-era documents ${\approx}0.37$.
We run 8~seeds (42, 99, 123, 200, 314, 500, 777, 1000) for each condition and keep the run with the highest final log-likelihood.

\subsection{Results}

Table~\ref{tab:validation} shows the percentage of each extracted topic's probability mass falling on each true topic's word group, for the best seed in each condition.

\begin{table}[ht]
\centering
\caption{Validation: $K_{\text{true}}=6$, $K_{\text{extract}}=6$, best of 8~seeds.
  Each row is an extracted topic; each column is the mass on a true topic's words.
  Bold entries mark dominant assignments ($\geq99.9\%$).}
\label{tab:validation}
\begin{tabular}{@{}l rrrrrr@{}}
\toprule
 & \multicolumn{3}{c}{\makebox[0pt]{\textbf{Sparse era (1850)}}} & \multicolumn{3}{c}{\makebox[0pt]{\textbf{Dense era (1950)}}} \\
\cmidrule(lr){2-4} \cmidrule(l){5-7}
\textbf{Topic} & \textbf{T1} & \textbf{T2} & \textbf{T3} & \textbf{T4} & \textbf{T5} & \textbf{T6} \\
 & \texttt{abc} & \texttt{def} & \texttt{ghi} & \texttt{ABC} & \texttt{DEF} & \texttt{GHI} \\
\midrule
\multicolumn{7}{@{}l}{\makebox[0pt][l]{\textit{Weighted + local $\boldsymbol\alpha$ (seed 500, LL $= -169{,}147$)}}} \\
T0 &  0.0 &  0.0 &  0.0 & \textbf{100.0} &  0.0 &  0.0 \\
T1 &  0.0 & \textbf{100.0} &  0.0 &  0.0 &  0.0 &  0.0 \\
T2 &  0.0 &  0.0 &  0.0 &  0.0 & \textbf{100.0} &  0.0 \\
T3 &  0.0 &  0.0 & \textbf{100.0} &  0.0 &  0.0 &  0.0 \\
T4 &  0.0 &  0.0 &  0.0 &  0.0 &  0.0 & \textbf{100.0} \\
T5 & \textbf{100.0} &  0.0 &  0.0 &  0.0 &  0.0 &  0.0 \\
\cmidrule{1-7}
\multicolumn{7}{@{}l@{}}{\makebox[0pt][l]{\textbf{6/6 true topics recovered (3 sparse + 3 dense)}}} \\
\midrule
\multicolumn{7}{@{}l}{\makebox[0pt][l]{\textit{Unweighted (seed 99, LL $= -170{,}514$)}}} \\
T0 &  0.0 &  0.0 &  0.0 &  0.0 & \textbf{100.0} &  0.0 \\
T1 &  0.0 &  0.0 &  0.0 & \textbf{100.0} &  0.0 &  0.0 \\
T2 &  0.0 &  0.0 &  0.0 &  0.0 &  0.0 & \textbf{100.0} \\
T3 &  0.0 &  0.0 &  0.0 &  0.0 & \textbf{100.0} &  0.0 \\
T4 &  0.0 & \textbf{99.9} &  0.1 &  0.0 &  0.0 &  0.0 \\
T5 & 60.2 &  0.0 & 39.8 &  0.0 &  0.0 &  0.0 \\
\cmidrule{1-7}
\multicolumn{7}{@{}l@{}}{\makebox[0pt][l]{\textbf{4/6 true topics recovered (1 sparse + 3 dense)}}} \\
\multicolumn{7}{@{}l@{}}{\makebox[0pt][l]{\textbf{DEF duplicated (T0, T3); abc+ghi collapsed into one blob}}} \\
\bottomrule
\end{tabular}
\end{table}

\subsection{Recovered $\boldsymbol\alpha$ Vectors}

The weighted + local model learns a $K$-vector $\boldsymbol\alpha(y)$ for each year.
The asymmetric Minka estimator drives inactive topics toward $\alpha_k \approx 0$ and active topics toward their true Dirichlet concentration:

\begin{center}
\begin{tabular}{@{}l rrrrrr r@{}}
\toprule
\textbf{Year} & $\alpha_0$ & $\alpha_1$ & $\alpha_2$ & $\alpha_3$ & $\alpha_4$ & $\alpha_5$ & $\alpha_\Sigma$ \\
\midrule
1850 & 0.000 & 0.107 & 0.000 & 0.115 & 0.000 & 0.120 & 0.342 \\
1950 & 0.302 & 0.000 & 0.270 & 0.000 & 0.282 & 0.000 & 0.854 \\
\bottomrule
\end{tabular}
\end{center}

At 1850, only topics T1, T3, T5 (the sparse-era topics, mapped to \texttt{def}, \texttt{ghi}, \texttt{abc} respectively) have nonzero $\alpha_k$, with a sum of~$0.342 \approx 3 \times 0.1$ (the true $\operatorname{Dir}(0.1,0.1,0.1)$).
At 1950, only topics T0, T2, T4 (the dense-era topics, mapped to \texttt{ABC}, \texttt{DEF}, \texttt{GHI}) have nonzero $\alpha_k$, with a sum of~$0.854 \approx 3 \times 0.3$ (the true $\operatorname{Dir}(0.3,0.3,0.3)$).
The estimator correctly identifies which topics are active in each era and recovers the per-era concentration to within 15\% of the true value.

\subsection{Analysis}

\paragraph{Weighted model.}
Every extracted topic places 100\% of its probability mass on a single true topic's three words.
All six true topics---including the three from the sparse era, which contributed only 5{,}000 of the 105{,}000 tokens---are recovered perfectly.
The density weights ($7.3\times$ for 1850, $0.37\times$ for 1950) equalise the per-era influence on the topic-word sufficient statistics, and the asymmetric $\boldsymbol\alpha$ estimator correctly learns the era-specific sparsity structure.
Seven of the eight seeds land in the same optimum (LL~$\approx -169{,}148$); only one seed finds a different basin.

\paragraph{Unweighted model.}
All three dense-era topics are recovered, but the model \emph{duplicates} one of them: extracted topics T0 and T3 both map to \texttt{DEF}.
One sparse-era topic (\texttt{def}) is recovered at 99.9\%, but the remaining two sparse-era topics (\texttt{abc} and \texttt{ghi}) are collapsed into a single blob (60/40\%).
Despite having the benefit of asymmetric $\alpha$ optimisation, the unweighted model cannot overcome the 20$\times$ token imbalance: it wastes one of its six topic slots duplicating a dense-era topic rather than resolving sparse-era structure.
The eight seeds split across two LL basins (3 seeds at ${\sim}{-}170{,}500$; 5 seeds at ${\sim}{-}173{,}060$), with the worse basin showing 0/3 sparse recovery.

\paragraph{Interpretation.}
Without weighting, the 100{,}000 dense-era tokens overwhelm the 5{,}000 sparse-era tokens in the topic-word counts.
Even with asymmetric $\alpha$ optimisation (which represents a strictly more powerful model than standard symmetric LDA), the unweighted sampler allocates its capacity to the dense era.
With weighting, each era's tokens carry balanced influence, and the model distributes its capacity proportionally to the number of true themes regardless of document frequency.

%% ===================================================================
\section{Underfitting Experiment ($K < K_{\text{true}}$)}
%% ===================================================================

We now ask what happens when the model is given fewer topics than truly exist.
This is the normal situation in practice: the analyst does not know $K_{\text{true}}$ and must choose~$K$.

\subsection{Setup}

The corpus is identical to Appendix~A (550~documents, 105{,}000 tokens, 6~true topics, heterogeneous $\operatorname{Dir}(0.1)$/$\operatorname{Dir}(0.3)$).
We fit $K=4$ topics, forcing the model to merge at least two pairs of true topics.
All other settings match Appendix~A: initial $\alpha=1$, initial $\beta=1$, up to 2{,}000 iterations, \texttt{--optimize-interval\;50}, \texttt{--converge\;1e-6}, best of 8~seeds.

\subsection{Results}

Table~\ref{tab:underfit} shows the topic-word mass breakdown with $K=4$.

\begin{table}[ht]
\centering
\caption{Underfitting: $K_{\text{true}}=6$, $K_{\text{extract}}=4$, best of 8~seeds.
  Each row is an extracted topic; each column is the mass on a true topic's words.}
\label{tab:underfit}
\begin{tabular}{@{}l rrrrrr@{}}
\toprule
 & \multicolumn{3}{c}{\makebox[0pt]{\textbf{Sparse era (1850)}}} & \multicolumn{3}{c@{}}{\makebox[0pt]{\textbf{Dense era (1950)}}} \\
\cmidrule(lr){2-4} \cmidrule(l){5-7}
\textbf{Topic} & \textbf{T1} & \textbf{T2} & \textbf{T3} & \textbf{T4} & \textbf{T5} & \textbf{T6} \\
 & \texttt{abc} & \texttt{def} & \texttt{ghi} & \texttt{ABC} & \texttt{DEF} & \texttt{GHI} \\
\midrule
\multicolumn{7}{@{}l}{\makebox[0pt][l]{\textit{Weighted + local $\boldsymbol\alpha$ (seed 200, LL $= -196{,}056$, converged iter 800)}}} \\
T0 &  0.0 & \textbf{100.0} &  0.0 &  0.0 &  0.0 &  0.0 \\
T1 &  0.0 &  0.0 &  0.0 & \textbf{100.0} &  0.0 &  0.0 \\
T2 &  0.0 &  0.0 &  0.0 &  0.0 & 48.9 & 51.1 \\
T3 & 60.2 &  0.0 & 39.8 &  0.0 &  0.0 &  0.0 \\
\cmidrule{1-7}
\multicolumn{7}{@{}l@{}}{\makebox[0pt][l]{\textbf{2 pure + 2 merged: balanced 2:2 split across eras}}} \\
\midrule
\multicolumn{7}{@{}l}{\makebox[0pt][l]{\textit{Unweighted (seed 314, LL $= -173{,}059$)}}} \\
T0 &  0.0 &  0.0 &  0.0 &  0.0 & \textbf{100.0} &  0.0 \\
T1 & 36.8 & 38.9 & 24.3 &  0.0 &  0.0 &  0.0 \\
T2 &  0.0 &  0.0 &  0.0 &  0.0 &  0.0 & \textbf{100.0} \\
T3 &  0.0 &  0.0 &  0.0 & \textbf{100.0} &  0.0 &  0.0 \\
\cmidrule{1-7}
\multicolumn{7}{@{}l@{}}{\makebox[0pt][l]{\textbf{0 sparse + 3 dense: 3:1 split}}} \\
\multicolumn{7}{@{}l@{}}{\makebox[0pt][l]{\textbf{Sparse era collapsed into one blob}}} \\
\bottomrule
\end{tabular}
\end{table}

\subsection{Recovered $\boldsymbol\alpha$ Vectors ($K=4$)}

With only 4~topics, the weighted model maps two topics to each era.
The asymmetric estimator correctly assigns nonzero $\alpha_k$ only to the topics active in each era:

\begin{center}
\begin{tabular}{@{}l rrrr r@{}}
\toprule
\textbf{Year} & $\alpha_0$ & $\alpha_1$ & $\alpha_2$ & $\alpha_3$ & $\alpha_\Sigma$ \\
\midrule
1850 & 0.085 & 0.000 & 0.000 & 0.155 & 0.240 \\
1950 & 0.000 & 0.300 & 0.543 & 0.000 & 0.843 \\
\bottomrule
\end{tabular}
\end{center}

At 1850, topics T0 (\texttt{def}) and T3 (\texttt{abc+ghi}) have nonzero $\alpha_k$, summing to $0.240$.
Since two of the three true sparse-era topics have been merged into T3, one ``slot'' of the true $\operatorname{Dir}(0.1,0.1,0.1)$ is absorbed, and the sum undershoots the true $3 \times 0.1 = 0.3$.
At 1950, topics T1 (\texttt{ABC}) and T2 (\texttt{DEF+GHI}) sum to $0.843 \approx 3 \times 0.3$, correctly recovering the dense era's total concentration.
The merged topic T2 absorbs two dense-era components, reflected in its larger $\alpha_2 = 0.543 \approx 2 \times 0.3$.

\subsection{Analysis}

With $K=4$ and $K_{\text{true}}=6$, the model must merge two pairs of true topics.
The question is \emph{which} merges occur.

\paragraph{Unweighted model.}
All three dense-era topics are recovered purely, and all three sparse-era topics are collapsed into a single blob (37/39/24\%).
The model devotes 3 of its 4 topics to the dense era and only 1 to the sparse era---a 3:1 split even though both eras have 3 true topics.
All 8~seeds converge to essentially the same configuration (LL range of only ${\sim}10$), confirming a single dominant basin.

\paragraph{Weighted model.}
The merging is distributed symmetrically across both eras:
\begin{itemize}
\item T0 recovers \texttt{def} purely (100\%).
\item T1 recovers \texttt{ABC} purely (100\%).
\item T3 merges two sparse-era topics (\texttt{abc}\;+\;\texttt{ghi}, 60/40).
\item T2 merges two dense-era topics (\texttt{DEF}\;+\;\texttt{GHI}, 49/51).
\end{itemize}
The model distributes its 4~topics in a 2:2 split between eras: each era gets one pure topic and one merged topic.
The weighted model converged at iteration~800.
The 8~seeds split across two basins---4 near LL~$\approx -197{,}000$ with the 2:2 pattern, 4 near LL~$\approx -226{,}600$ with a different merge---confirming that multi-seed selection is important for underfit models.

\paragraph{Interpretation.}
When the model is underfitting ($K < K_{\text{true}}$), it must make compromises.
The unweighted model \emph{always} sacrifices the sparse era first because dense-era tokens dominate the likelihood.
The weighted model spreads the compromise symmetrically: each era keeps one pure topic and surrenders one to merging.
This is exactly the behaviour predicted by the uniform-effective-mass property (Section~5): each era exerts equal influence, so the model splits its budget proportionally rather than starving the underrepresented period.

\end{document}
